\chapter{Objetivos}
\label{cap:objetivos}
% --------------------------------------------------------------------------

\section{Objetivo geral}
\label{sec:objetivo-geral}

Desenvolver, validar e qualificar computacionalmente um arcabouço autônomo de otimização volumétrica e geométrica de metamateriais auxéticos aplicados ao chassi primário de nanossatélites CubeSat 1U em liga aeroespacial AlSi10Mg fabricado por Fusão em Leito de Pó a Laser (L-PBF), integrando Simulação em Elementos Finitos de Alta Fidelidade (FEA), Modelos Substitutos Neurais Guiados por Física (\textit{Physics-Guided Machine Learning}) e Grafos de Fluxo Determinísticos de Engenharia de Sistemas, com vistas a maximizar a atenuação elastodinâmica passiva sob o envelope normativo da norma NASA GSFC-STD-7000A (GEVS) e minimizar a massa estrutural, garantindo a integridade estática de escoamento e a vida em fadiga estocástica de voo.

\section{Objetivos específicos}
\label{sec:objetivos-especificos}

Para viabilizar o cumprimento sistemático do objetivo geral e fornecer comprovação quantitativa para as três hipóteses científicas enunciadas no Capítulo~\ref{cap:introducao}, a pesquisa desdobra-se em nove objetivos específicos interdependentes:

\begin{alineas}
  \item \textbf{Parametrização geométrica e preservação cinemática de contato (Objetivo 1):} Parametrizar a célula unitária auxética reentrante tridimensional e os painéis estruturais do chassi 1U em função do vetor de variáveis de projeto $\mathbf{x} = [\theta, t, l, h]^\top$, preservando a continuidade monolítica dos quatro trilhos-guia longitudinais de \rail{} com cantos arredondados ($R \ge \SI{1.0}{\milli\metre}$) e rugosidade superficial $Ra \le \SI{1.6}{\micro\metre}$, garantindo o correto acoplamento cinemático e eliminando o risco de engripamento mecânico (\textit{jamming}) ou soldagem a frio no interior do dispensador orbital P-POD;
  \item \textbf{Formalização das restrições de manufatura aditiva metálica DfAM (Objetivo 2):} Mapear, programar e validar os filtros determinísticos de Projeto para Manufatura Aditiva (\textit{Design for Additive Manufacturing} -- DfAM) específicos para o processo L-PBF em liga AlSi10Mg, estabelecendo como critérios mandatórios de viabilidade: ângulo de sobre-salha autossuportado $\theta_{\text{overhang}} \ge \SI{35.0}{\degree}$, espessura mínima de costela $t \ge \SI{0.60}{\milli\metre}$ para estabilidade da poça de fusão, espaçamento livre entre hastes $\ge \SI{1.50}{\milli\metre}$ e orifícios de drenagem de pó não fundido com diâmetro desobstruído $d_{\text{drain}} \ge \SI{2.0}{\milli\metre}$;
  \item \textbf{Planejamento e amostragem do espaço de projeto viável (Objetivo 3):} Estruturar uma campanha de amostragem por amostragem quase-Monte Carlo via Hipercubo Latino (\textit{Latin Hypercube Sampling} -- LHS) com 250 pontos no espaço viável $[\theta, t, l, h]^\top$, gerando a base de dados de treinamento e validação para o modelo substituto por meio do modelo multiescala celular calibrado por elementos finitos;
  \item \textbf{Automação da esteira elastodinâmica em código aberto (Objetivo 4):} Desenvolver scripts em Python fortemente tipados para a automação da esteira de simulação em elementos finitos de alta fidelidade em ambiente aberto (Code\_Aster, CalculiX e Gmsh Tet10 em Linux/WSL), executando análises modais de autovalores via algoritmo de Sorensen/Lanczos (\code{CALC_MODES}) e análises espectrais de vibração aleatória (\code{DYNA_ALEA_MODAL}) sob excitação PSD de \SI{14.1}{\Grms} (\SIrange{20}{2000}{\hertz}) da norma NASA GSFC-STD-7000A;
  \item \textbf{Formulação e treinamento da rede neural guiada por física (Objetivo 5):} Projetar, treinar e validar uma arquitetura de rede neural profunda com blocos residuais (\textit{Physics-Guided ResNet}), empregando função de perda composta ($\mathcal{L}_{\text{total}}$) que incorpora termos fenomenológicos de penalização assintótica de Gibson-Ashby e penalidades de violação de monotonicidade elástica ($\partial f_1 / \partial t > 0$), garantindo predições fisicamente consistentes de submilissegundos ($R^2 > 0{,}98$);
  \item \textbf{Implementação do grafo determinístico de governança multiagente (Objetivo 6):} Instanciar uma máquina de estados finitos determinística baseada na biblioteca LangGraph (\code{StateGraph}), governada por contratos de dados imutáveis fortemente tipados (\code{CubeSatState}), segregando papéis especializados de portão DfAM, inferência substituta neural e qualificação normativa espacial, com geração de trilha de auditoria completa com carimbo temporal UTC e justificativas formais de descarte;
  \item \textbf{Otimização evolutiva multiobjetivo e caracterização de Pareto (Objetivo 7):} Acoplar o algoritmo genético evolutivo NSGA-II à esteira de inferência neural para mapear a fronteira de Pareto de compromisso entre minimização de massa estrutural e minimização de transmissibilidade dinâmica na carga útil, submetido a restrições ativas de rigidez ($f_1 \ge \SI{100.0}{\hertz}$), limite elástico de escoamento ($MS_{\text{yield}} > 0$) e acúmulo de dano por fadiga estocástica de Steinberg ($D \le \num{0.25}$);
  \item \textbf{Tomada de decisão multicritério e validação cruzada de alta ordem (Objetivo 8):} Selecionar o candidato ótimo de voo pelo método de decisão multicritério TOPSIS, submetendo-o a validação cruzada independente de alta fidelidade no solver Code\_Aster (\textit{Ground Truth}) e CalculiX, acompanhado de estudo rigoroso de convergência de malha de segunda ordem (Gmsh Tet10) via Índice de Convergência de Malha (GCI de Roache);
  \item \textbf{Avaliação de integridade em voo e benchmarking aeroespacial (Objetivo 9):} Executar a qualificação dinâmica estrutural completa do chassi auxético ótimo frente aos requisitos de lançamento — incluindo a simulação de dinâmica explícita não linear de choque no dispensador P-POD via OpenRadioss —, e confrontar os resultados de desempenho com o chassi convencional usinado maciço em liga de Al 6061-T6, comprovando quantitativamente a superioridade em alívio de massa e atenuação vibratória.
\end{alineas}

\section{Articulação metodológica, fases operacionais e matriz de rastreabilidade}
\label{sec:articulacao-objetivos}

A execução dos nove objetivos específicos é estruturada em quatro fases operacionais consecutivas. A Tabela~\ref{tab:matriz-objetivos} apresenta a matriz de rastreabilidade, correlacionando cada fase e objetivo específico aos métodos aplicados, critérios de aceitação e entregáveis mensuráveis da pesquisa.

\begin{table}[htb]
  \caption{Matriz de rastreabilidade, métodos computacionais e critérios de aceitação dos objetivos específicos}
  \label{tab:matriz-objetivos}
  \small
  \SingleSpacing
  \begin{tabularx}{\textwidth}{@{}>{\raggedright\arraybackslash}p{2.2cm} >{\raggedright\arraybackslash}p{2.6cm} >{\raggedright\arraybackslash}p{3.8cm} >{\raggedright\arraybackslash}X@{}}
    \toprule
    \textbf{Fase operacional} & \textbf{Objetivo específico} & \textbf{Método / Ferramenta} & \textbf{Critério de aceitação / Métrica} \\
    \midrule
    \textbf{Fase I:} \newline Fundamentação Geométrica e DfAM & 
    Obj. 1: Parametrização B-Rep \newline
    Obj. 2: Filtros DfAM L-PBF & 
    OpenCASCADE / CadQuery / Python B-Rep; equações cinemáticas de célula reentrante. & 
    Quatro trilhos maciços contínuos de \rail{}; $\theta_{\text{overhang}} \ge \SI{35}{\degree}$; $t \ge \SI{0.60}{\milli\metre}$; furos de pó $\phi \ge \SI{2.0}{\milli\metre}$. \\
    \midrule
    \textbf{Fase II:} \newline Síntese Numérica CAE e DoE & 
    Obj. 3: Amostragem LHS \newline
    Obj. 4: Solvers FEA abertos & 
    Latin Hypercube Sampling (250 pts); Gmsh Tet10; Code\_Aster (\code{CALC_MODES}, \code{DYNA_ALEA_MODAL}). & 
    Geração do banco de dados multiescala calibrado; extração de $f_1$, $\sigma_{3\sigma}$ e transmissibilidade sob PSD de \SI{14.1}{\Grms}. \\
    \midrule
    \textbf{Fase III:} \newline Metamodelo e Governança & 
    Obj. 5: Physics-Guided ResNet \newline
    Obj. 6: Grafo LangGraph & 
    PyTorch ResNet com perdas fenomenológicas Gibson-Ashby e monotonicidade; LangGraph StateGraph determinístico. & 
    Acurácia de regressão $R^2 > 0{,}98$ com erro RMSE $< \SI{3}{\percent}$; tempo de inferência $< \SI{1}{\milli\second}$; trilha de auditoria UTC sem falha. \\
    \midrule
    \textbf{Fase IV:} \newline Otimização, Validação e Benchmarking & 
    Obj. 7: NSGA-II \newline
    Obj. 8: TOPSIS e GCI \newline
    Obj. 9: OpenRadioss e Benchmark & 
    Algoritmo evolutivo NSGA-II; TOPSIS; convergência GCI (Roache); verificação CalculiX; dinâmica explícita OpenRadioss. & 
    $f_1 \ge \SI{100.0}{\hertz}$; massa estrutural $\le \SI{0.150}{\kilo\gram}$ (alívio $> \SI{50}{\percent}$); atenuação $> \SI{10}{\decibel}$; $MS_{\text{yield}} > 0$; dano Steinberg $D \le \num{0.25}$; GCI $< \SI{1.5}{\percent}$. \\
    \bottomrule
  \end{tabularx}
  \fonte{Elaborada pelo autor (2026).}
\end{table}

Essa estrutura garante que cada entrega técnica atue como pré-requisito rigorosamente validado para o avanço da esteira, assegurando que o chassi de voo resultante atinja conformidade espacial inquestionável perante as normas da NASA e da AEB.
